\chapter{The Black Cat}
Before the cat came, of course there had to be a broomstick.
\par
In the bazaar of one of our hill stations is an old junk shop-dirty, dingy and dark-in which I often potter about looking for old books or Victorian bric-a-brac. Sometimes, one comes across useful household items, but I do not usually notice these. I was, however, attracted to an old but well-preserved broom standing in a corner of a shop A long-handled broom was the rooms of my cottage, and I did not enjoy bending over double when using the common short-handled \textit{jharoo}.
\par
The old broom was priced at ten rupees. I haggled with the shopkeeper and got it for five.
\par
It was a strong broom, full of character, and I used it to good effect almost every morning. And there this story might have ended-or would never have begun-if I had not found the large black cat sitting on the garden wall.
\par
The black cat had bright yellow eyes, and it gave me a long penetrating look, as through it were summing up my possibilities as an exploitable human. Through it miaowed once or twice, I paid it no attention. I did not care much for cats. But when I went indoors, I found that the cat had followed me and begun scratching at the pantry door.
\par
It must be hungry, I thought, and gave it some milk. The cat lapped up the milk, purring deeply all the while, then sprang up on a cupboard and made itself comfortable.
\par
Well, for several days, there was no getting rid of the cat. It seemed completely at home, and merely tolerated my presence in the house. It was more interested in my broom than me, and would dance and skitter around the broom whenever I was sweeping the rooms. And when the broom was resting against the wall, the cat would slide up to it, rubbing itself against the handle and purring loudly.
\par
A cat and a broomstick-the combination was suggestive, full of possibilities . . . The cottage was old, almost a hundred years old, and I wondered about the kind of tenants it might have had during these long years. I had been in the cottage for only a year. And through it stood alone in the midst of a forest of Himalayan oaks, I had never encountered any ghosts or spirits.
\par
Miss Bellows came to see me in the middle of July. I heard the tapping of a walking stick on the rocky path outside the cottage, a tapping that stopped near the gate.
\par
'Mr Bond!' called an imperious voice. 'Are you home?'
\par
I had to been doing some gardening; I looked up to find an elderly straight-backed Englishwoman peering at me over the gate.
\par
'Good evening' I said, dropping my hoe.
\par
'I believe you have my cat,' said Miss Bellows.
\par
Through I had not met the lady before, I knew her by name and reputation. She was the oldest resident of the hill station.
'I do have a cat,' I said, 'through it's probably more correct to say that the cat has me. If she's your cat, you're welcome to her. Why don't you come in while I looked for her?'
\par
Mr Bellows stepped in. She wore a rather old-fashioned black dress, and her ancient but strong walnut stick had two or three curves in it and a knob instead of a handle.
\par
She made herself comfortable in an armchair while I went in search of the cat. But the cat was on one of her mysterious absences, and through I called for her in my most persuasive manner, she did not respond. I knew she was probably quite near. But cats are like that-preserve obstinate creatures.
\par
When the finally I returned to the sitting room, there was the cat, curled up on Miss Bellows' lap.
\par
'Well, you've got her, I see. Would you like some tea before you go?'
\par
'No, thank you,' said Miss Bellows. 'I don't drink tea.'
\par
'Something stronger, perhaps. A little brandy?' She looked up at me rather sharply. Disconcerted, I hastened to add, 'Not that I drink much, you know. I keep a little in the house for emergencies. It helps ward off colds and things. It's particularly good for-er-well, for colds,' I finished lamely.
\par
'I see your kettle's boiling,' she said. 'Can I have some hot water?'
\par
'Hot water? Certainly.' I was a little puzzled, but I did not want to antagonize Miss Bellows at our first meeting.
\par
'Thank you. And a glass.'
\par
She took a glass and I went to get the kettle, From the pocket of her voluminous dress, she extracted two small packets, similar to those containing chemists' powders. Opening both packets, she poured first a purple powder and then a crimson powder into the glass. Nothing happened.
\par
'Now the water, please,' she said.
\par
'It's boiling hot!'
\par
'Never mind.'
\par
I poured boiling water into her glass, and there was a terrific fizzing and bubbling as the frothy stuff rose to the rim. It gave off a horrible stench. The potion was so hot that I thought it would crack the glass, but before this could happen, Miss Bellows put it to her lips and drained the contents.
\par
'I think I'll be going now,' she said, putting the glass down and smacking her lips. The cat, tail in the air, voiced its agreement. Said Miss Bellows, 'I'm much obliged to you, young man.'
\par
'Don't mention it,' I said humbly. 'Always at your service.'
\par
She gave me her thin body hand and held mine in an icy grip.
\par
I saw Miss Bellows and the black cat to the gate and returned pensively to my sitting room. Living alone was beginning to tell on my nerves and imagination. I made a half-hearted attempt to laugh at my fancies, but the laugh stuck in my throat. I couldn't help noticing that the broom was missing from its corner.
\par
I dashed out of the cottage and looked up and down the path. There was no one to be seen. In the gathering darkness, I could hear Miss Bellows' laughter, followed by a snatch of song:
\vspace{0.4in}
\par
\textit{With the darkness round me growing,}
\par
\textit{And the moon behind my hat,}
\par
\textit{You will soon have trouble knowing}
\par
\textit{Which is witch and witch's cat.}
\vspace{0.4in}
\par
Something whirred overhead like a Diwali rocket.
\par
I looked up and saw them silhouetted against the rising moon. Miss bellows and her cat were riding away on my broomstick.
